% 原书印刷页1；PDF物理页32.章标题由主文件给出.

\sourcepageanchor{32}

\begin{quote}
例如，你永远无法预言某一个人会做什么，却可以准确地说出一群人平均而言会怎样行动.个体各不相同，但百分比保持不变.统计学家就是这样说的.
\begin{flushright}夏洛克·福尔摩斯\\《四签名》\end{flushright}
\end{quote}
概率论（probability theory）是一切统计学赖以建立的基础，它提供了对总体、试验，以及几乎任何可以视为随机现象的事物进行建模的方法.借助这些模型，统计学家能够仅考察整体的一部分，就对总体作出推断.

概率论有着悠久而丰富的历史，至少可以追溯到十七世纪.当时，应朋友德·梅雷骑士之请，帕斯卡和费马用数学形式描述了赌博中的胜算.

本章并不打算全面介绍概率论；在如此有限的篇幅内作这样的尝试未免鲁莽.我们的目标是概述概率论中一些对学习统计学至关重要的基本思想.

正如统计学建立在概率论的基础之上，概率论又建立在集合论的基础之上，因此我们从集合论开始.
